\chapter{Google OAuth Login}

Email/password login works, but many users expect a "Sign in with Google"
button too -- one less password to create and remember. This chapter walks
through wiring Google OAuth into the Task Management System so that, after
Google verifies who the user is, your own backend still issues the same JWT
cookie used everywhere else in the app. Nothing downstream (the
\texttt{protect} middleware, \texttt{authorize}, controllers) needs to know
or care whether a user logged in with a password or with Google.

\section{The Big Picture}

The key idea: your React app never handles the user's Google password, and
your Express server never talks to Google's login page directly. Instead,
Google's own SDK on the frontend handles the login UI and hands your backend
a signed \textbf{ID token} proving the user's identity, which your backend
verifies with Google before trusting it.

\begin{diagrambox}[Google OAuth Flow]
\begin{verbatim}
 1. React renders Google's "Sign in with Google" button
        |  user clicks, Google popup/redirect handles login
        v
 2. Google -> Authorization (user approves access)
        |
        v
 3. Callback -> React receives an ID token from Google
        |
        v
 4. React -> POST /api/auth/google  { idToken }
        |
        v
 5. Express -> verifies token with Google
             (OAuth2Client.verifyIdToken)
        |
        v
 6. Express -> Find user by email, or Create if first login
        |
        v
 7. Express -> Generate JWT + set cookie (same as normal login)
        |
        v
 8. React -> receives success response, user is authenticated
\end{verbatim}
\end{diagrambox}

\section{Google Cloud Console Setup}

Before any code, register the app in the Google Cloud Console (APIs \&
Services $\rightarrow$ Credentials $\rightarrow$ Create OAuth Client ID):

\begin{itemize}
  \item Choose \textbf{Web application} as the application type.
  \item Add your frontend origin (e.g. \texttt{http://localhost:5173}) under
    \textbf{Authorized JavaScript origins}.
  \item Add your callback/redirect URI under \textbf{Authorized redirect
    URIs}, if using the redirect-based flow rather than the popup-based
    Google Identity Services button.
  \item Copy the generated \textbf{Client ID} and \textbf{Client Secret}.
\end{itemize}

\begin{lstlisting}[language=JavaScript]
# .env
GOOGLE_CLIENT_ID=1234567890-abc123.apps.googleusercontent.com
GOOGLE_CLIENT_SECRET=GOCSPX-xxxxxxxxxxxxxxxxxxxxxxxx
\end{lstlisting}

\begin{warnbox}[Common Error: redirect\_uri\_mismatch]
This is the single most common Google OAuth setup error. Google requires the
redirect URI your app actually sends to \emph{exactly} match one of the URIs
registered in the Cloud Console -- byte for byte. Common mismatches:
\begin{itemize}
  \item \texttt{http://localhost:5173} registered, but the app runs on
    \texttt{http://localhost:5173/} (trailing slash) or a different port.
  \item \texttt{http} registered in development, but \texttt{https} used in
    production (these are two separate URIs to Google, not the same one).
  \item A subdomain difference, e.g. \texttt{app.example.com} registered but
    \texttt{www.app.example.com} used.
\end{itemize}
Fix it by copying the exact URI your app sends (visible in the error message
or network tab) and adding that precise string to Authorized redirect URIs
in the Cloud Console, or by correcting your app's configured redirect URI to
match what's registered.
\end{warnbox}

\section{Backend: Verifying the ID Token}

Using the \texttt{google-auth-library} package, the backend receives the ID
token from the frontend and verifies it directly with Google's servers --
this confirms the token is genuine and was issued for \emph{your} Client ID.

\begin{lstlisting}[language=JavaScript]
// controllers/authController.js
import { OAuth2Client } from "google-auth-library";
import User from "../models/User.js";
import { generateToken } from "../utils/generateToken.js";

const client = new OAuth2Client(process.env.GOOGLE_CLIENT_ID);

export const googleLogin = async (req, res, next) => {
  try {
    const { idToken } = req.body;

    const ticket = await client.verifyIdToken({
      idToken,
      audience: process.env.GOOGLE_CLIENT_ID,
    });

    const payload = ticket.getPayload();
    const { email, name, picture } = payload;

    let user = await User.findOne({ email });

    if (!user) {
      // First-time Google login: create the account.
      // No password is set for OAuth-only accounts.
      user = await User.create({
        name,
        email,
        avatar: picture,
        password: undefined,
        role: "user",
      });
    }

    const token = generateToken(user._id);

    res.cookie("token", token, {
      httpOnly: true,
      secure: process.env.NODE_ENV === "production",
      sameSite: "lax",
      maxAge: 7 * 24 * 60 * 60 * 1000,
    });

    res.status(200).json({
      success: true,
      data: { id: user._id, name: user.name, email: user.email, role: user.role },
    });
  } catch (err) {
    next(err);
  }
};
\end{lstlisting}

\begin{lstlisting}[language=JavaScript]
// routes/authRoutes.js
router.post("/auth/google", googleLogin);
\end{lstlisting}

\section{Handling the Password Field}

Because Google-authenticated users never set a password on your system, the
User schema's \texttt{password} field must be optional for this to work:

\begin{lstlisting}[language=JavaScript]
// models/User.js
password: { type: String, required: false },
\end{lstlisting}

\begin{notebox}[NOTE]
If a user later tries to log in with email/password using an account that
was created via Google (and therefore has no password hash), the login
controller's \texttt{bcrypt.compare} call needs a guard: check
\texttt{if (!user.password)} first and return a clear "This account uses
Google Sign-In" message instead of attempting to compare against
\texttt{undefined}.
\end{notebox}

\section{Subsequent Logins}

On every later login, the same \texttt{googleLogin} controller runs:
\texttt{User.findOne(\{ email \})} finds the existing account (skipping the
creation branch), and a fresh JWT cookie is issued exactly as it was the
first time. From the rest of the application's perspective -- middleware,
controllers, the frontend's authenticated state -- a Google-authenticated
session and a password-authenticated session are indistinguishable; both end
with the same signed JWT cookie.

\begin{tipbox}[TIP]
Keep the Google verification logic isolated to this one controller. Every
other part of the backend (\texttt{protect}, \texttt{authorize}, task
controllers) only ever deals with your own JWT, never with Google tokens
directly. This is what makes adding a second OAuth provider later
(GitHub, Facebook) a matter of adding another controller, not rewriting
existing auth logic.
\end{tipbox}

\whatwhyhow
{An endpoint that accepts a Google-issued ID token from the frontend, verifies it server-side with Google, finds or creates a matching user, and issues your app's own JWT cookie.}
{Letting the frontend alone decide "this user is authenticated with Google" would be trivially spoofable; server-side verification against Google's servers is the only trustworthy point, and re-using your existing JWT flow keeps the rest of the app provider-agnostic.}
{Verify with \texttt{OAuth2Client.verifyIdToken}, \texttt{findOrCreate} the user by email, then call the exact same \texttt{generateToken} + \texttt{res.cookie} steps used for normal login.}
